% --- MODULE 8 (PART 2): Problems (BST, Tries, Graphs) ---
% --- This file is now UPDATED with problems from GFG PDFs ---

\section{Module 8 (Part 2): Problems (Post-BST)}

\subsection{Module 6: Trees, BSTs, Tries}

\subsubsection{Problem: Validate Binary Search Tree (Blind 75)}
\begin{itemize}
\item \textbf{Problem:} Determine if a binary tree is a valid BST.
\item \textbf{Solution Idea:} Use a recursive DFS. Pass down the valid range (min\_val, max\_val) allowed for the current node. The invariant must hold for all ancestors.
\item isValid(node, min, max):
\item if node.key <= min || node.key >= max: return false.
\item return isValid(node.left, min, node.key) \&\& isValid(node.right, node.key, max).
\item \textbf{Complexity:} $O(n)$ time, $O(h)$ space.
\end{itemize}

\subsubsection{Problem: BST Deletion (from GFG Binary-Search-Tree.pdf)}
\begin{itemize}
\item \textbf{Problem:} Delete a node z from a BST (Lec 25).
\item \textbf{Solution Idea:}
\begin{itemize}
\item \textbf{Case 1 (0 or 1 child):} Replace z with its child.
\item \textbf{Case 2 (2 children):} Find z's \textbf{inorder successor} y (the minimum node in z.right). y is guaranteed to have at most one (right) child.
\item Copy y.key into z.key.
\item Recursively delete y (which is now a Case 1 problem) from the right subtree.
\end{itemize}
\item \textbf{Complexity:} $O(h)$ time.
\end{itemize}

\subsubsection{Problem: Kth Smallest Element in a BST (Blind 75)}
\begin{itemize}
\item \textbf{Problem:} Find the $k$-th smallest element in a BST.
\item \textbf{Solution Idea 1 (Inorder):} Perform an inorder traversal. Stop at the $k$-th visited node. $O(n)$ time.
\item \textbf{Solution Idea 2 (Augmented BST):} (Lec 26) If nodes store size(left\_subtree), you can find it in $O(h)$ time.
\item \textbf{Complexity 2:} $O(h)$ time.
\end{itemize}

\subsubsection{Problem: Lowest Common Ancestor of a BST (from GFG)}
\begin{itemize}
\item \textbf{Problem:} Find the LCA of two nodes p and q in a BST.
\item \textbf{Solution Idea:} This is much simpler than a general BT. Start from the root.
\item if p.key < root.key \&\& q.key < root.key: recurse on root.left.
\item if p.key > root.key \&\& q.key > root.key: recurse on root.right.
\item else: root is the LCA (the point where p and q split).
\item \textbf{Complexity:} $O(h)$ time, $O(1)$ space (iterative).
\end{itemize}

\subsubsection{Problem: Convert Sorted Array to Balanced BST (from GFG)}
\begin{itemize}
\item \textbf{Problem:} Given a sorted array, create a height-balanced BST.
\item \textbf{Solution Idea:} This is a Divide and Conquer approach.
\begin{enumerate}
\item The \textbf{middle} element of the array arr[mid] becomes the root.
\item The left subarray arr[0...mid-1] recursively forms the root.left subtree.
\item The right subarray arr[mid+1...n-1] recursively forms the root.right subtree.
\end{enumerate}
\item \textbf{Complexity:} $O(n)$ time (visits each node once).
\end{itemize}

\subsubsection{Problem: Lowest Common Ancestor of a Binary Tree (Blind 75)}
\begin{itemize}
\item \textbf{Problem:} Find the LCA of two nodes $p$ and $q$ in a general binary tree.
\item \textbf{Solution Idea:} Use a recursive DFS find(node, p, q).
\item \textbf{Base Case:} if node == NULL or node == p or node == q, return node.
\item \textbf{Recurse:} left = find(node.left), right = find(node.right).
\item \textbf{Combine:} if left \&\& right: return node (it's the LCA). else: return left or right (whichever is not NULL).
\item \textbf{Complexity:} $O(n)$ time, $O(h)$ space.
\end{itemize}

\subsubsection{Problem: Construct Tree from Preorder and Inorder (Blind 75)}
\begin{itemize}
\item \textbf{Problem:} Given preorder and inorder traversal arrays, construct the binary tree.
\item \textbf{Solution Idea:} (1) First element of preorder is the root. (2) Find root in inorder to split it into left/right subarrays. (3) Use lengths of these subarrays to split preorder. (4) Recurse.
\item \textbf{Key:} Use a hash map for inorder indices to get $O(1)$ lookup.
\item \textbf{Complexity:} $O(n)$ time, $O(n)$ space.
\end{itemize}

\subsubsection{Problem: Binary Tree Maximum Path Sum (Hard, Blind 75)}
\begin{itemize}
\item \textbf{Problem:} Find the maximum sum of any path in a binary tree.
\item \textbf{Solution Idea:} Use a recursive DFS max\_gain(node). It returns the max path sum \textit{extendable upwards}. It also updates a global max\_sum with the max path sum \textit{contained at this node} (node.key + left\_gain + right\_gain).
\item \textbf{Complexity:} $O(n)$ time, $O(h)$ space.
\end{itemize}

\subsubsection{Problem: AVL Tree Insertion (from GFG Advanced-Data-Structure.pdf)}
\begin{itemize}
\item \textbf{Problem:} Insert a node into an AVL tree and rebalance (Lec 27-29).
\item \textbf{Solution Idea:} (1) Perform standard BST insert. (2) Walk back up from the new node to the root. (3) At each ancestor z, update its height. (4) Find the \textbf{first} ancestor z that violates the AVL property ($|h\_L - h\_R| > 1$). (5) Perform the correct \textbf{Rotation} (LL, RR, LR, RL) at z.
\item \textbf{Key:} Only \textbf{one} rotation (single or double) is needed to fix an insertion.
\item \textbf{Complexity:} $O(\log n)$ time.
\end{itemize}

\subsubsection{Problem: B-Tree Insertion (from GFG Advanced-Data-Structure.pdf)}
\begin{itemize}
\item \textbf{Problem:} Insert a key into a B-Tree (Lec 32).
\item \textbf{Solution Idea (Pre-emptive Split):} We \textit{never} insert into a full node.
\begin{enumerate}
\item Start at the root. As we traverse \textit{down} to find the correct leaf:
\item \textbf{If we encounter a full node x:} We \textbf{split x first} (promoting its median key to its parent) before moving down into the correct (now non-full) child.
\item \textbf{Finally:} We reach the target leaf, which is \textit{guaranteed} not to be full. Insert key $k$.
\end{enumerate}
\item \textbf{Complexity:} $O(\log\_t n)$ disk accesses.
\end{itemize}

\subsubsection{Problem: Implement Trie (Prefix Tree) (Blind 75)}
\begin{itemize}
\item \textbf{Problem:} Implement a Trie with insert, search, and startsWith.
\item \textbf{Solution Idea:} Use a TrieNode class with children[26] and isEndOfWord.
\item \textbf{Complexity:} All ops are $O(L)$ time, where $L$ is length of the word.
\end{itemize}

\subsubsection{Problem: Auto-complete feature using Trie (from GFG Advanced-Data-Structure.pdf)}
\begin{itemize}
\item \textbf{Problem:} Given a query string s, return all words in the dictionary that have s as a prefix.
\item \textbf{Solution Idea:} (1) search the Trie for the node corresponding to the prefix s. (2) If this node p exists, run a DFS from p to find all words in the subtree rooted at p.
\item \textbf{Complexity:} $O(L + K \cdot M)$ where $L$=length of prefix, $K$=number of suggestions, $M$=avg word length.
\end{itemize}

\subsection{Module 7: Graph Algorithms}

\subsubsection{Problem: Number of Islands (Blind 75)}
\begin{itemize}
\item \textbf{Problem:} Given a 2D grid of '1's (land) and '0's (water), count the islands.
\item \textbf{Solution Idea:} Iterate the grid. If grid[r][c] == '1' and unvisited: increment island\_count, and start a \textbf{BFS or DFS} to visit and mark all connected '1's.
\item \textbf{Complexity:} $O(N \cdot M)$ time.
\end{itemize}

\subsubsection{Problem: Topological Sort (from GFG Graphs.pdf)}
\begin{itemize}
\item \textbf{Problem:} Find a linear ordering of vertices in a DAG such that for every edge $(u, v)$, $u$ comes before $v$ (Lec 38).
\item \textbf{Solution Idea (DFS):} (1) Run DFS on the graph. (2) As each vertex \textbf{finishes} (is marked BLACK), insert it onto the \textbf{front} of a linked list. (3) This list is the topological sort.
\item \textbf{Correctness:} For an edge $(u,v)$, $v$ must finish before $u$ (otherwise it's a back edge/cycle). Since $u$ finishes later, it gets added to the front of the list, placing it before $v$.
\item \textbf{Complexity:} $O(n+m)$ time.
\end{itemize}

\subsubsection{Problem: Course Schedule (Lec 38, Blind 75)}
\begin{itemize}
\item \textbf{Problem:} Can you finish all courses given prerequisites?
\item \textbf{Solution Idea:} This is \textbf{Cycle Detection in a Directed Graph}.
\item \textbf{Algorithm (DFS):} Use 3 "colors": WHITE (unvisited), GRAY (in current recursion stack), BLACK (finished). If DFS visits a GRAY node, a \textbf{back edge} (cycle) exists.
\item \textbf{Complexity:} $O(n+m)$ time.
\end{itemize}

\subsubsection{Problem: Detect Cycle in Undirected Graph (from GFG Graphs.pdf)}
\begin{itemize}
\item \textbf{Problem:} Check if an undirected graph has a cycle.
\item \textbf{Solution Idea 1 (DFS):} Run DFS. Keep track of each node's parent. If you find an edge to an already-visited node v that is \textbf{not} your parent, you have found a back edge/cycle.
\item \textbf{Solution Idea 2 (DSU):} (Lec 41) Iterate all edges $(u, v)$. For each edge, if Find(u) == Find(v): a cycle exists (they are already connected). else: Union(u, v).
\item \textbf{Complexity:} $O(n+m)$ for DFS. $O(m \cdot \alpha(n))$ for DSU.
\end{itemize}

\subsubsection{Problem: Snake and Ladder (from GFG Graphs.pdf)}
\begin{itemize}
\item \textbf{Problem:} Find the minimum number of dice throws to win a Snake and Ladder game.
\item \textbf{Solution Idea:} This is a \textbf{Shortest Path in an Unweighted Graph} problem.
\item \textbf{Graph:} The board (squares 1-100) is the set of $n$ vertices.
\item \textbf{Edges:} From each square i, there are 6 edges to i+1, ..., i+6 (dice rolls).
\item \textbf{Snakes/Ladders:} These are not edges, but "teleports". If a ladder is at i to j, you only have an edge to j, not i. (Model this in the adj. list).
\item \textbf{Algorithm:} Run \textbf{BFS} from square 1. The dist[100] is the answer.
\item \textbf{Complexity:} $O(n)$ time (where $n=100$).
\end{itemize}

\subsubsection{Problem: Word Search (Backtracking/DFS)}
\begin{itemize}
\item \textbf{Problem:} Find if a word exists in a 2D grid. Cannot reuse a cell.
\item \textbf{Solution Idea:} \textbf{DFS with backtracking}. Start DFS from any cell matching word[0]. Mark cell as visited, explore 4 neighbors, then un-mark cell.
\item \textbf{Complexity:} $O(N \cdot M \cdot 3^L)$ where $L$ is word length.
\end{itemize}

\subsubsection{Problem: Kruskal's Algorithm for MST (from GFG Graphs.pdf)}
\begin{itemize}
\item \textbf{Problem:} Find a Minimum Spanning Tree (Lec 41).
\item \textbf{Solution Idea:} (1) \textbf{Sort} all $m$ edges by weight ($O(m \log m)$). (2) Initialize a \textbf{DSU} with $n$ sets. (3) Iterate sorted edges $(u, v)$: if Find(u) != Find(v): add $(u, v)$ to MST and Union(u, v).
\item \textbf{Complexity:} $O(m \log m)$ or $O(m \log n)$ (dominated by sort).
\end{itemize}